\chapter*{Preliminary Philosophical Considerations}


\section*{Introduction}

A peculiar intellectual roadblock halted my initial attempt to conduct this research even before it began.
The fundamental question I faced concerned the nature of research itself.
How could I proceed to conduct research without first explicating what research precisely is?
Upon inquiry, I came to regard research as the process of attaining the aim of any scientific endeavor:
\textit{expanding the body of human knowledge}.

This definition, however, generated a vexing new question and led me down the rabbit hole of epistemology:
\textit{What is knowledge}?
To answer this question — or any question — I would need to \textit{know} the answer, thus employing knowledge to define the concept of knowledge, which is evidently tautological.

The reader might possess a peripheral awareness of the problem I faced and perhaps has conquered this obstacle in their own work in some way or another.
I, however, could not proceed without addressing this foundational issue.

The root cause of these problems lies in the central question of epistemology:
\textit{What is knowledge}?
Linguistically, this question constitutes a sensible inquiry.
Yet, none of the words in this dissertation would hold any value without establishing a foundational epistemic framework in which to house them.
Thu, this preliminary chapter aims to construct this philosophical bedrock and lay the mandatory groundwork for conducting any further work.


\section*{Motivation}

This dissertation is a scientific endeavor.
The aim of any scientific endeavor is the expansion of knowledge.
To expand knowledge, we must first define what knowledge is—that is, what can and what cannot be known.

To pose a research question, we must first distinguish the space of sensible answers from nonsensical and meaningless ones. To configure and execute a research method, we must operate within an epistemic framework that justifies our approach. Without a clear understanding of what constitutes knowledge, any attempt to contribute to the existing body of knowledge would be unfounded.

I will not delve explicitly into the different epistemological positions, nor will I provide an exhaustive overview, as it would be superfluous for the purposes of this dissertation. Instead, I will argue my case through which my epistemological positioning will become apparent.

I acknowledge that, in the format of this text, I can use mere words to communicate my thoughts. My first request, therefore, is for the reader to conspire with me and accept that what I will attempt to convey can, in fact, be verbalized. Otherwise, how may the nature of knowledge be examined when the words employed to that end are required to be true and thus contain knowledge themselves?

\textit{Note}: I consciously chose the term \textit{word} to differentiate it from \textit{language} and other symbols. See also the \enquote{Triangle of Reference} \parencite{ogden1923meaning}.

\section*{The Postulate of Shared Meaning}

To begin, let us confront a meta-ontological question and ask how we can define (in words) any concept—including knowledge—that we can grasp with our mind. Axiomatically, we must accept that our cognition maps the concepts in our mind \textit{appropriately} onto the words we use. This statement remains a necessary assumption for any meaningful discussion. Any attempt to investigate what constitutes this appropriateness (i.e., circumstances, time, individualities) is futile because the spoken words are the only gateway to the mind of the speaker \parencite{wittgenstein1953philosophical}.

Any attempt to bypass words and focus directly on our own mind would inevitably lead to methodological solipsism \parencite[136]{putnam1975meaning}, rendering any communication about the concepts contained within the mind redundant. As Putnam illustrates in his \enquote{Brain in a Vat} thought experiment, without a connection to external reality, our words lose their meaning \parencite{putnam1981reason}.

It follows, then, that words can only obtain \textit{meaning} when their sense—not their truth—can be validated against something outside the mind of the speaker. That is, when no one hears me, speaking itself becomes nonsensical \parencite{wittgenstein1922tractatus}. I can call the sky \enquote{blue} today, \enquote{green} tomorrow, and \enquote{red} the day after. Each statement would be true to me and would have to pass as knowledge because there is no way (and no need) to prove it. Under these circumstances, everything I say becomes true by default.

We have arrived at my first epistemic stance: the \textit{Postulate of Shared Meaning}, which states that knowledge can only exist intersubjectively. This means that the words I use to verbalize my mind's concepts must correspond—when interpreted—to \textit{equivalent} concepts in the mind of the listener. But how can we be sure of that?

\section*{The Postulate of Pragmatism}

Even when no other human is listening to us, we are never truly speaking into the void. Nature is always present. We cannot deceive nature nor escape it. However, nature does not communicate in human language. Thus, we cannot ask it to explain itself verbally. We may only pose questions in the form of actions and observe the responses. Over time, we can build theories about nature and express them in words, under the assumption that other humans interpreting these words share the same understanding and use the same terminology to describe the phenomena observed in nature \parencite[1011b]{aristotle2009metaphysics}. This forms the foundation of research in the natural sciences.

The static and tangible subject matter in the natural sciences allows us to construct theories and adopt methods of fallibilism. The natural sciences circumvent the need for a strict definition of knowledge by grounding their arguments in pragmatism. But can we apply the same methodological approach and epistemological paradigms when considering a contingent subject like the human mind \parencite{frank2007introduction}? The answer is complex.

So, again, how can we be certain that the words we speak will evoke the same mental concepts in the listener as in ourselves? To avoid endless debates about the \enquote{nature of reality}, let us adopt a phenomenological approach to this question \parencite{husserl2002ideas}. Putnam's \enquote{Twin Earth} thought experiment highlights the challenges of semantic externalism and the meaning of words \parencite{putnam1975meaning}.

One approach is to consider words as our portal into the human mind (as discussed earlier) and engage in a dialogue with others about what constitutes knowledge. But which words—which knowledge—will ultimately be accepted? We tend to accept that which is most useful and effective in facilitating successful actions in practical contexts. This means that truth is not an abstract, static concept but is determined by the practical outcomes of holding certain beliefs. This aligns with the \enquote{Pragmatic Theory of Truth} \parencite{james1907pragmatism,capps2023pragmatic}.

This leads to my second epistemic assertion: the \textit{Postulate of Pragmatism}, which states that knowledge must lead to justified action. Knowledge is not contingent upon external validation or empirical evidence but on its inherent capacity to produce practical outcomes. We may not dissect knowledge into further components like justification and belief. As Timothy Williamson argues in \textit{Knowledge and Its Limits}, knowledge is a mental state that cannot be analyzed into simpler components \parencite{williamson2002knowledge}.

But what is a \enquote{problem}, and what constitutes a \enquote{successful action}? Epistemology reaches its limits here and answers with \textit{belief}. To delve deeper, we would need to explore existentialist philosophy to find answers, which I will briefly do. The question is, how can we construct a bridge to close the gap between the pragmatism required for knowledge and the evaluation of the proof of pragmatic ideas when we have no mirror—like nature—against which to check our knowledge?

\section*{The Postulate of the Will}

When knowledge can only be expressed in words—as we have established—it appears that only under the assumption that a limited number of word configurations exist can we speak of \enquote{contributions to the body of knowledge}. However, this assumption does not hold. First, even the countable set of words that constitute our vocabulary grows, and there is no inherent limit to this growth—a priori or empirically. Second, the set of combinations of words—their configurations—is infinite.

We could attempt to bypass this by accepting that our mind itself is limited in its capacity to create mental images out of these word configurations and that, eventually, clusters of images emerge that are \enquote{similar enough} to speak of equivalency within these clusters. This, however, requires a similarity metric—not necessarily quantifiable.

But how do we determine this similarity? By advancing one's will, as posited by philosophers such as Arthur Schopenhauer, Friedrich Nietzsche, and Jean-Paul Sartre. In this sense, as the will is continuously evolving and at any point there exists one thinkable yet unattainable optimal solution to advancing it, knowledge is at any point limited but requires extension when considered on a continuum.

This leads me to my final epistemic position and my confession to following a perspectivist approach: the \textit{Postulate of the Will}, which states that it is our individual will that interprets and gives meaning to our experiences and actions \parencite{schopenhauer2009world}. Schopenhauer's concept of the will as the fundamental reality suggests that our desires and motivations shape our perceptions and understandings.

To conclude, I can only attempt to communicate my own \textit{perspective} and strive to argue, to the best of my ability, how pragmatically useful my thoughts are—for me. I cannot—nor will I try—to convince the reader of anything outside the scope of my perspective.

\section*{Summary}

I sought to define what knowledge is. I first approached this through the philosophy of language, questioning how we can define, using words, \textit{anything}. We established that shared meaning is a prerequisite for meaningful communication.

We then examined how we can ensure that we actually mean the same thing by the words we use. By borrowing methods from the natural sciences, I posited that it is the pragmatic content that imbues knowledge with value. This means that knowledge is not something that statements possess inherently but rather that knowledge is what leads to the most beneficial outcomes in our actions. It also implies that we can never know anything with absolute certainty; we can only know \textit{enough} for a specific \textit{goal}.

To advance these thoughts further and answer what pragmatism is and by what metric it can be evaluated, I turned away from blind faith and looked into the philosophy of existentialism, employing the concept of the \textit{will} as propagated by prominent philosophers. Will is highly individualistic and cannot—nor should it—be further investigated in this context.

In adopting this perspectivist approach, I recognize that even if the ideas I present have been articulated before, the way I express them—through my own words and perspective—is inherently original. The scope of generalization is thus both personal and universal; while my insights stem from my individual will, they may resonate with others due to shared human experiences.

In conclusion, I must adopt a perspectivist approach in this dissertation and convey how things appear to \textit{me} and what insights \textit{I} need to advance \textit{my} will. Fear not, though. As we are all humans and my readers are presumably scholars like myself, it is fair to assume that my insights will stem from an appropriately similar will and thus exhibit enough congruency to be useful for my readers.

These considerations are particularly pertinent in the field of information systems, where the utility and applicability of knowledge are paramount. Understanding that knowledge is shaped by individual perspectives and wills allows us to design systems that are more adaptable and responsive to the needs of their users.